% Options for packages loaded elsewhere
\PassOptionsToPackage{unicode}{hyperref}
\PassOptionsToPackage{hyphens}{url}
%
\documentclass[
]{article}
\usepackage{amsmath,amssymb}
\usepackage{lmodern}
\usepackage{iftex}
\ifPDFTeX
  \usepackage[T1]{fontenc}
  \usepackage[utf8]{inputenc}
  \usepackage{textcomp} % provide euro and other symbols
\else % if luatex or xetex
  \usepackage{unicode-math}
  \defaultfontfeatures{Scale=MatchLowercase}
  \defaultfontfeatures[\rmfamily]{Ligatures=TeX,Scale=1}
\fi
% Use upquote if available, for straight quotes in verbatim environments
\IfFileExists{upquote.sty}{\usepackage{upquote}}{}
\IfFileExists{microtype.sty}{% use microtype if available
  \usepackage[]{microtype}
  \UseMicrotypeSet[protrusion]{basicmath} % disable protrusion for tt fonts
}{}
\makeatletter
\@ifundefined{KOMAClassName}{% if non-KOMA class
  \IfFileExists{parskip.sty}{%
    \usepackage{parskip}
  }{% else
    \setlength{\parindent}{0pt}
    \setlength{\parskip}{6pt plus 2pt minus 1pt}}
}{% if KOMA class
  \KOMAoptions{parskip=half}}
\makeatother
\usepackage{xcolor}
\setlength{\emergencystretch}{3em} % prevent overfull lines
\providecommand{\tightlist}{%
  \setlength{\itemsep}{0pt}\setlength{\parskip}{0pt}}
\setcounter{secnumdepth}{-\maxdimen} % remove section numbering
\ifLuaTeX
  \usepackage{selnolig}  % disable illegal ligatures
\fi
\IfFileExists{bookmark.sty}{\usepackage{bookmark}}{\usepackage{hyperref}}
\IfFileExists{xurl.sty}{\usepackage{xurl}}{} % add URL line breaks if available
\urlstyle{same} % disable monospaced font for URLs
\hypersetup{
  hidelinks,
  pdfcreator={LaTeX via pandoc}}

\author{}
\date{}

\begin{document}

\documentclass{article}
\usepackage{amsmath}
\usepackage{amsfonts}

\begin{document}

\section*{Lecture Scribe}

\subsection*{CSE400 – Fundamentals of Probability in Computing}

\subsection*{Lecture L15\_16\_17\_18: Joint Random Variables + Joint Distributions + Examples}

---

\section*{Outline}

The lecture outline lists the following topics:

\begin{itemize}
\item Why Study Joint Random Variables?
\item Discrete Case
\begin{itemize}
\item Joint PMF Table Representation
\item Example: Rolling Two Dice
\end{itemize}
\item Continuous Case
\begin{itemize}
\item Geometric Interpretation
\item Worked Example
\end{itemize}
\item Joint Cumulative Distribution Function
\begin{itemize}
\item Definition of Joint CDF
\item Continuous Case
\item Discrete Case
\item Properties of the Joint CDF
\end{itemize}
\item Comprehensive Example
\end{itemize}

---

\section*{1. Why Study Joint Random Variables?}

In science and in real life, we are often interested in \textbf{two (or more) random variables at the same time}.

\subsection*{Examples}

\begin{itemize}
\item Height and weight of giraffes
\item IQ and birthweight of children
\item Frequency of exercise and rate of heart disease
\item Air pollution level and respiratory illness rate
\item Number of Facebook friends and age of members
\end{itemize}

\subsection*{Example \#1: Smart Transportation}

In a smart traffic system, two random variables are defined:

\begin{itemize}
\item $(X =)$ Vehicle Speed
\item $(Y =)$ Traffic Density
\end{itemize}

Question stated:

\begin{itemize}
\item Can speed be high when the traffic density is also high?
\end{itemize}

Answer shown:

\begin{itemize}
\item Yes / No
\end{itemize}

Statement given:

\begin{itemize}
\item So, $(X)$ and $(Y)$ are correlated random variables.
\end{itemize}

Real ITS Application (probabilistic analysis):

\begin{itemize}
\item Adaptive traffic signals
\item Congestion prediction
\item Autonomous vehicle routing
\end{itemize}

\subsection*{Example \#2: Wireless Network Performance}

In a wireless system, two important random variables are:

\begin{itemize}
\item $(X =)$ Signal-to-Noise Ratio (SNR)
\item $(Y =)$ Data Throughput
\end{itemize}

Why Joint Random Variables?

\begin{itemize}
\item Throughput depends on SNR
\item Both vary due to fading, interference, and congestion
\end{itemize}

Question stated:

\begin{itemize}
\item If SNR is high, will throughput always be high?
\end{itemize}

Answer shown:

\begin{itemize}
\item Yes / No
\end{itemize}

Reasoning stated:

\begin{itemize}
\item Network Congestion (condition relationships)
\end{itemize}

Network Applications:

\begin{itemize}
\item 5G scheduling
\item adaptive modulation
\item network planning
\end{itemize}

\subsection*{Example \#3}

Weather Prediction:

\begin{itemize}
\item $(X =)$ Rainfall
\item $(Y =)$ Temperature
\end{itemize}

Application:

\begin{itemize}
\item weather forecasting
\end{itemize}

\subsection*{Example \#4}

Drone Delivery:

\begin{itemize}
\item $(X =)$ Wind Speed
\item $(Y =)$ Delivery Time
\end{itemize}

Applications:

\begin{itemize}
\item logistics
\item UAV systems
\end{itemize}

\subsection*{Example \#5}

\begin{itemize}
\item Roll two dice $((X \text{ and } Y))$ using phones
\item $((X + Y = 7 \text{ and } X > Y))$
\item This requires joint PMF
\end{itemize}

\subsection*{Statements concluding the section}

In such situations, the random variables have a \textbf{joint distribution} that allows us to compute probabilities of events involving both variables and understand the \textbf{relationship} between the variables.

This is simplest when the variables are \textbf{independent}.

When they are not, we use \textbf{covariance} and \textbf{correlation} as measures of the nature of the dependence between them.

---

\section*{2. Discrete Case – Joint PMF}

\subsection*{Setup}

Suppose $(X)$ and $(Y)$ are discrete random variables.

\[
X \in \{x_1,x_2,\ldots,x_n\}, \qquad
Y \in \{y_1,y_2,\ldots,y_m\}
\]

The ordered pair $((X,Y))$ takes values in the product set:

\[
\{(x_1,y_1),(x_1,y_2),\ldots,(x_n,y_m)\}
\]

\subsection*{Definition (Joint PMF)}

\[
p(x_i,y_j)=P(X=x_i,;Y=y_j)
\]

This gives the probability of the joint outcome occurring simultaneously.

---

\section*{3. Joint PMF Table Representation}

The joint PMF is organized in a table. The table has rows indexed by values of $(X)$ and columns indexed by values of $(Y)$, with entries $(p(x_i,y_j))$.

\subsection*{Key Properties}

\[
0 \le p(x_i,y_j) \le 1
\]

\[
\sum_{i=1}^{n}\sum_{j=1}^{m} p(x_i,y_j)=1
\]

---

\section*{4. Example: Rolling Two Dice}

\subsection*{Experiment}

Roll two fair dice.

\subsection*{Let}

\begin{itemize}
\item $(X =)$ value on the first die
\item $(T =)$ total (sum) of both dice
\end{itemize}

Since each outcome $((i,j))$ has probability $\left(\frac{1}{36}\right)$, we construct the joint probability table of $((X,T))$.

\subsection*{Joint Probability Table of $((X,T))$}

\[
\begin{array}{c|ccccccccccc}
X \backslash T & 2 & 3 & 4 & 5 & 6 & 7 & 8 & 9 & 10 & 11 & 12 \\ \hline
1 & \frac{1}{36} & \frac{1}{36} & \frac{1}{36} & \frac{1}{36} & \frac{1}{36} & \frac{1}{36} & 0 & 0 & 0 & 0 & 0 \\
2 & 0 & \frac{1}{36} & \frac{1}{36} & \frac{1}{36} & \frac{1}{36} & \frac{1}{36} & \frac{1}{36} & 0 & 0 & 0 & 0 \\
3 & 0 & 0 & \frac{1}{36} & \frac{1}{36} & \frac{1}{36} & \frac{1}{36} & \frac{1}{36} & \frac{1}{36} & 0 & 0 & 0 \\
4 & 0 & 0 & 0 & \frac{1}{36} & \frac{1}{36} & \frac{1}{36} & \frac{1}{36} & \frac{1}{36} & \frac{1}{36} & 0 & 0 \\
5 & 0 & 0 & 0 & 0 & \frac{1}{36} & \frac{1}{36} & \frac{1}{36} & \frac{1}{36} & \frac{1}{36} & \frac{1}{36} & 0 \\
6 & 0 & 0 & 0 & 0 & 0 & \frac{1}{36} & \frac{1}{36} & \frac{1}{36} & \frac{1}{36} & \frac{1}{36} & \frac{1}{36}
\end{array}
\]

\subsection*{Observations}

\begin{itemize}
\item Entries are either $(0)$ or $\left(\frac{1}{36}\right)$
\item $(X)$ and $(T)$ are \textbf{not independent}
\end{itemize}

---

\section*{5. Continuous Case – Joint PDF}

The lecture states that the continuous case is analogous to the discrete case:

\begin{itemize}
\item Discrete values $(\rightarrow)$ Continuous intervals
\item Joint PMF $(\rightarrow)$ Joint PDF
\item Sums $(\rightarrow)$ Double integrals
\end{itemize}

If

\[
X \in [a,b], \qquad Y \in [c,d],
\]

then the pair $((X,Y))$ takes values in the product region:

\[
[a,b]\times[c,d]
\]

\subsection*{Definition (Joint PDF)}

A function $(f(x,y))$ is called the \textbf{joint probability density function} of $(X)$ and $(Y)$ if

\[
P\big((X,Y)\text{ lies in a small rectangle of area }dx,dy\big)\approx f(x,y),dx,dy
\]

\subsection*{Geometric Interpretation}

The region shown is the rectangle $([a,b]\times[c,d])$, and inside it a small rectangle of dimensions $(dx)$ and $(dy)$ is marked with probability

\[
\text{Prob.} = f(x,y),dx,dy
\]

\subsection*{Key Properties}

\[
f(x,y)\ge 0
\]

\[
\int_c^d \int_a^b f(x,y),dx,dy = 1
\]

Additional statement given:

\begin{itemize}
\item The joint PDF $(f(x,y))$ may take values greater than $(1)$.
\item It is a probability density, not a probability.
\end{itemize}

---

\section*{6. Worked Example}

\subsection*{Given}

\[
f(x,y)=4xy, \qquad 0 \le x \le 1,; 0 \le y \le 1
\]

\subsection*{Tasks}

\begin{enumerate}
\item Show that $(f(x,y))$ is a valid joint PDF.
\item Visualize the event
\[
A=\{X<0.5,;Y>0.5\}
\]
\item Compute $(P(A))$
\end{enumerate}

Jointly, $((X,Y))$ takes values in the \textbf{unit square}

\[
[0,1]^2
\]

The figure shows the region $(A)$ inside the unit square corresponding to $(X<0.5)$ and $(Y>0.5)$.

\subsection*{Example Solution}

\subsubsection*{Step 1: Check Validity (Normalization)}

\[
\int_0^1 \int_0^1 4xy,dx,dy
\]

\[
===========================
\]

\[
\# \int_0^1 \left(\int_0^1 4xy,dx\right)dy
\]

\[
\# \int_0^1 2y,dy
\]

\[
[y^2]_0^1
=1
\]

Hence $(f(x,y)\ge 0)$ and total probability is $(1)$. So it is a \textbf{valid joint PDF}.

\subsubsection*{Step 2: Compute $(P(A))$}

\[
P(A)=\int_0^{0.5}\int_{0.5}^{1} 4xy,dy,dx
\]

\[
=\int_0^{0.5}[2xy^2]_{0.5}^{1},dx
\]

\[
=\int_0^{0.5}2x(1-0.25),dx
\]

\[
=\int_0^{0.5}\frac{3}{2}x,dx
\]

\[
=\frac{3}{16}
\]

---

\section*{7. Joint Cumulative Distribution Function (Joint CDF)}

Suppose $(X)$ and $(Y)$ are jointly distributed random variables.

We use the notation:

\[
X \le x,; Y \le y
\]

to mean the event:

\[
\{X \le x \text{ and } Y \le y\}
\]

\subsection*{Definition}

The joint cumulative distribution function is defined as

\[
F(x,y)=P(X \le x,; Y \le y)
\]

---

\section*{8. Joint CDF – Continuous Case}

If $(X)$ and $(Y)$ are continuous random variables with joint density $(f(x,y))$ over $([a,b]\times[c,d])$, then the joint CDF is obtained by integrating the density:

\[
F(x,y)=\int_c^y \int_a^x f(u,v),du,dv
\]

\subsection*{Recovering the Joint PDF}

Since there are two variables, we use partial derivatives:

\[
f(x,y)=\frac{\partial^2}{\partial x \partial y}F(x,y)
\]

---

\section*{9. Joint CDF – Discrete Case}

If $(X)$ and $(Y)$ are discrete random variables with joint PMF $(p(x_i,y_j))$, then the joint CDF is obtained by summing probabilities:

\[
F(x,y)=\sum_{x_i\le x}\sum_{y_j\le y} p(x_i,y_j)
\]

\subsection*{Interpretation}

\begin{itemize}
\item Add all probabilities in the table
\item For rows with $(x_i \le x)$
\item And columns with $(y_j \le y)$
\end{itemize}

---

\section*{10. Properties of the Joint CDF}

The joint CDF $(F(x,y))$ of $(X)$ and $(Y)$ must satisfy several properties.

\subsection*{Property 1}

\[
F(x,y)\text{ is non-decreasing.}
\]

If $(x)$ or $(y)$ increase, then $(F(x,y))$ must stay constant or increase.

\subsection*{Property 2}

\[
F(x,y)=0 \text{ at the lower-left of the joint range.}
\]

If the lower-left corner is

\[
(-\infty,-\infty),
\]

then

\[
\lim_{(x,y)\to(-\infty,-\infty)}F(x,y)=0.
\]

\subsection*{Property 3}

\[
F(x,y)=1 \text{ at the upper-right of the joint range.}
\]

If the upper-right corner is

\[
(\infty,\infty),
\]

then

\[
\lim_{(x,y)\to(\infty,\infty)}F(x,y)=1.
\]

---

\section*{11. End Slide}

The final slide contains:

\[
\text{Thank You!}
\]

\end{document}

\end{document}
